\documentclass[UTF8,fontset=fandol,zihao=5]{ctexart}

\usepackage[a4paper,margin=2cm]{geometry}
\usepackage{amsmath,amssymb,mathtools,bm,esint}
\usepackage{booktabs,tabularx,array,longtable,multirow}
\usepackage{xcolor}
\usepackage[most]{tcolorbox}
\usepackage{tikz}
\usetikzlibrary{
  positioning,
  calc,
  decorations.pathreplacing,
  arrows.meta
}

\usepackage{enumitem}
\usepackage{indentfirst}
\usepackage{fancyhdr}
\usepackage{hyperref}

% relationbox 中知识关系图的统一样式
\tikzset{
  relnode/.style={
    font=\small,
    align=center
  },
  relline/.style={
    thick
  },
  relarrow/.style={
    -{Stealth},
    thick
  }
}

\DeclareTextFontCommand{\textbf}{\heiti\color{main}}

\definecolor{main}{HTML}{8E2C3A}
\definecolor{mainlight}{HTML}{FFFFFF}
\definecolor{method}{HTML}{8E5D64}
\definecolor{methodlight}{HTML}{F8EFF1}

\hypersetup{colorlinks=true,linkcolor=main,urlcolor=main}
\setlength{\parindent}{2em}
\linespread{1.5}
\setlength{\abovedisplayskip}{5pt}
\setlength{\belowdisplayskip}{5pt}
\setlength{\abovedisplayshortskip}{3pt}
\setlength{\belowdisplayshortskip}{3pt}
\setlist[itemize]{leftmargin=2em,itemsep=0.1em,topsep=0.2em}
\setlist[enumerate]{leftmargin=2.3em,itemsep=0.1em,topsep=0.2em}

\ctexset{
  section={format=\centering\heiti\zihao{3}\color{main},
           name={第,章},number=\chinese{section},aftername=\quad,
           beforeskip=0.9em,afterskip=0.5em},
  subsection={format=\heiti\zihao{5}\color{main},
              name={,、},number=\chinese{subsection},aftername=,
              beforeskip=0.6em,afterskip=0.22em},
  subsubsection={format=\heiti\zihao{5}\color{main},
                 number={（\chinese{subsubsection}）},aftername=,
                 beforeskip=0.55em,afterskip=0.15em}
}

\newtcolorbox{relationbox}{
  enhanced,breakable,colback=mainlight,colframe=main,
  boxrule=0.65pt,arc=1.4mm,left=2mm,right=2mm,top=1.2mm,bottom=1.2mm
}
\newtcolorbox{methodbox}[1][]{
  enhanced,breakable,colback=methodlight,colframe=method,
  boxrule=0.65pt,arc=1.4mm,left=2mm,right=2mm,top=1.2mm,bottom=1.2mm,
  title=#1,fonttitle=\kaishu,fontupper=\kaishu,
  before upper={\setlength{\parindent}{2em}}
}

\newtcolorbox{method}[1][]{
  enhanced,breakable,colback=methodlight,colframe=method,
  boxrule=0.65pt,arc=1.4mm,left=2mm,right=2mm,top=1.2mm,bottom=1.2mm,
  title={#1},fonttitle=\heiti\bfseries
}

\newcommand{\norm}[1]{\left\lVert#1\right\rVert}
\newcommand{\Prj}{\operatorname{Prj}}

\pagestyle{fancy}
\fancyhf{}
\fancyhead[C]{\small\color{method}解析几何笔记}
\fancyfoot[C]{\thepage}
\renewcommand{\headrulewidth}{0.35pt}

\begin{document}

\begin{center}
  {\heiti\zihao{1}\color{main} 解析几何笔记}
\end{center}

\subsection{平面}

设平面 $\Pi$ 法向量为 $\bm n=(A,B,C)\ne\bm 0$，$M(x_0,y_0,z_0)\in\Pi$。

\begin{description}[style=multiline,leftmargin=7em,labelwidth=6em,itemsep=0.65em]
  \item[点法式方程]
  $A(x-x_0)+B(y-y_0)+C(z-z_0)=0,\quad A^2+B^2+C^2\ne0$

  \item[一般式方程]
  $Ax+By+Cz+D=0,\quad A^2+B^2+C^2\ne0$

  \hfill（$D=-Ax_0-By_0-Cz_0$）

  \item[截距式方程]
  $\displaystyle\frac{x}{a}+\frac{y}{b}+\frac{z}{c}=1,\quad abc\ne0$

  \hfill（$a,b,c$ 分别为平面与 $x,y,z$ 三轴的截距。）

  \item[三点式方程]
  $\displaystyle
  \begin{vmatrix}
    x-x_1 & y-y_1 & z-z_1\\
    x-x_2 & y-y_2 & z-z_2\\
    x-x_3 & y-y_3 & z-z_3
  \end{vmatrix}=0$
\end{description}

\begin{method}[【注意】与谁平行谁为零]
如与 $x$ 轴平行则 $x=0$，与 $yOz$ 平面平行则 $y=z=0$；

过原点 $D=0$。
\end{method}

\begin{description}[style=multiline,leftmargin=7em,labelwidth=6em]
  \item[平面束方程]
  $\lambda(A_1x+B_1y+C_1z+D_1)+\mu(A_2x+B_2y+C_2z+D_2)=0,$

  \hfill $\lambda^2+\mu^2\ne0$
\end{description}

\subsection{直线}

设直线 $l$ 方向向量 $\bm s=(l,m,n)\ne\bm 0$，$M(x_0,y_0,z_0)\in l$。

\begin{description}[style=multiline,leftmargin=7em,labelwidth=6em,itemsep=0.7em]
  \item[参数方程]
  $\displaystyle
  \begin{cases}
    x=x_0+lt,\\
    y=y_0+mt,\\
    z=z_0+nt,
  \end{cases}\quad t\text{ 为参数}$

  \item[标准方程]
  $\displaystyle\frac{x-x_0}{l}=\frac{y-y_0}{m}=\frac{z-z_0}{n}$

  \hfill（又称点向式方程或对称方程）

  \item[一般方程]
  $\displaystyle
  \begin{cases}
    A_1x+B_1y+C_1z+D_1=0,\quad A_1^2+B_1^2+C_1^2\ne0,\\
    A_2x+B_2y+C_2z+D_2=0,\quad A_2^2+B_2^2+C_2^2\ne0
  \end{cases}$

  \hfill（又称普通方程）

  \item[两点式方程]
  $\displaystyle\frac{x-x_1}{x_2-x_1}=\frac{y-y_1}{y_2-y_1}=\frac{z-z_1}{z_2-z_1}$
\end{description}

\subsection{夹角}

\begin{align*}
\text{面面夹角：}\quad
\cos\langle\bm n_1,\bm n_2\rangle
&=\frac{|\bm n_1\cdot\bm n_2|}{\norm{\bm n_1}\norm{\bm n_2}},\\[0.4em]
\text{线面夹角：}\quad
\sin\langle\bm n_1,\bm s_1\rangle
&=\frac{|\bm n_1\cdot\bm s_1|}{\norm{\bm n_1}\norm{\bm s_1}},\\[0.4em]
\text{线线夹角：}\quad
\cos\langle\bm s_1,\bm s_2\rangle
&=\frac{|\bm s_1\cdot\bm s_2|}{\norm{\bm s_1}\norm{\bm s_2}}.
\end{align*}

\subsection{位置关系}

\subsubsection{平面}

\begin{align*}
\text{垂直：}\quad
\Pi_1\perp\Pi_2
&\Longleftrightarrow \bm n_1\perp\bm n_2
\Longleftrightarrow \bm n_1\cdot\bm n_2=0
\Longleftrightarrow A_1A_2+B_1B_2+C_1C_2=0,\\[0.6em]
\text{平行：}\quad
\Pi_1\parallel\Pi_2
&\Longleftrightarrow
\frac{A_1}{A_2}=\frac{B_1}{B_2}=\frac{C_1}{C_2}\ne\frac{D_1}{D_2}
\Longrightarrow \bm n_1\parallel\bm n_2
\Longleftrightarrow \bm n_1\times\bm n_2=\bm 0,\\[0.6em]
\text{重合：}\quad
\Pi_1\text{ 与 }\Pi_2\text{ 重合}
&\Longleftrightarrow
\frac{A_1}{A_2}=\frac{B_1}{B_2}=\frac{C_1}{C_2}=\frac{D_1}{D_2}
\Longrightarrow \bm n_1\parallel\bm n_2
\Longleftrightarrow \bm n_1\times\bm n_2=\bm 0.
\end{align*}

\subsubsection{直线}

\begin{align*}
\text{异面：}\quad
l_1\text{ 与 }l_2\text{ 异面}
&\Longleftrightarrow (\overrightarrow{M_1M_2},\bm s_1,\bm s_2)\ne0,\\
\text{共面：}\quad
l_1\text{ 与 }l_2\text{ 共面}
&\Longleftrightarrow (\overrightarrow{M_1M_2},\bm s_1,\bm s_2)=0.
\end{align*}

平行、垂直同上。

\subsection{距离}

点到平面距离：
\[
d=\Prj(\overrightarrow{P_0P_1},\bm n)
=\frac{|Ax_0+By_0+Cz_0+D|}{\sqrt{A^2+B^2+C^2}}.
\]

点到直线距离：
\[
d=\sqrt{\left|\overrightarrow{P_0P_1}\right|^2-
\left(\Prj(\overrightarrow{P_0P_1},\bm n)\right)^2}.
\]

\end{document}
